\documentclass[11pt]{ctexart}
\usepackage{amsmath,amssymb,amsthm}
\usepackage[margin=2.2cm,top=2.4cm,bottom=2.2cm]{geometry}
\usepackage{enumitem}
\usepackage{fancyhdr}
\usepackage{titlesec}
\usepackage{needspace}
\usepackage[hidelinks]{hyperref}
\setCJKmainfont{SimSun}[BoldFont=SimHei,ItalicFont=KaiTi]
\setCJKsansfont{SimHei}
\setCJKmonofont{FangSong}
\linespread{1.25}
\setlength{\parindent}{0pt}
\setlength{\parskip}{0.5em}
\allowdisplaybreaks
\pagestyle{fancy}
\fancyhf{}
\renewcommand{\headrulewidth}{0.4pt}
\fancyhead[L]{\small 高中数学联合竞赛一试·学案}
\fancyfoot[C]{\small\thepage}
\titleformat{\section}{\Large\bfseries}{}{0pt}{}
\titlespacing*{\section}{0pt}{0pt}{0.8em}
\newcommand{\bigproblem}[1]{\clearpage\section*{#1}}
\begin{document}
\sloppy
\newcommand{\blank}{\underline{\hspace{3em}}}
\clearpage
\section*{2013年全国高中数学联合竞赛一试（A卷）}

一、填空题：本大题共8小题，每小题8分，满分64分。

1. 设集合 $A=\{2,0,1,3\}$，集合 $B=\{x\mid -x\in A,\ 2-x^2\notin A\}$，则集合 $B$ 中所有元素的和为\blank{}。

2. 在平面直角坐标系 $xOy$ 中，点 $A,B$ 在抛物线 $y^2=4x$ 上，满足 $\overrightarrow{OA}\cdot\overrightarrow{OB}=-4$，$F$ 是抛物线的焦点，则 $S_{\triangle OFA}\cdot S_{\triangle OFB}=$\blank{}。

3. 在 $\triangle ABC$ 中，已知 $\sin A=10\sin B\sin C,\ \cos A=10\cos B\cos C$，则 $\tan A$ 的值为\blank{}。

4. 已知正三棱锥 $P-ABC$ 的底面边长为 $1$，高为 $\sqrt2$，则其内切球的半径为\blank{}。

5. 设 $a,b$ 为实数，函数 $f(x)=ax+b$ 满足：对任意 $x\in[0,1]$，有 $|f(x)|\le1$，则 $ab$ 的最大值为\blank{}。

6. 从 $1,2,\dots,20$ 中任取 $5$ 个不同的数，其中至少有 $2$ 个是相邻数的概率为\blank{}。

7. 若实数 $x,y$ 满足 $x-4\sqrt y=2\sqrt{x-y}$，则 $x$ 的取值范围为\blank{}。

8. 已知数列 $\{a_n\}$ 共有 $9$ 项，其中 $a_1=a_9=1$，且对每个 $i\in\{1,2,\dots,8\}$，均有 $\dfrac{a_{i+1}}{a_i}\in\left\{2,1,-\dfrac12\right\}$，则这样的数列的个数为\blank{}。

二、解答题：本大题共3小题，满分56分。


\clearpage
\textbf{9.（本题满分16分）}设正项数列 $\{x_n\}$ 满足 $S_n\ge 2S_{n-1},\ n\ge2$，这里 $S_n$ 是数列 $\{x_n\}$ 的前 $n$ 项和。求证：存在常数 $C>0$，使得 $x_n\ge C\cdot2^n$ 对任意正整数 $n$ 成立。

\clearpage
\textbf{10.（本题满分20分）}在平面直角坐标系 $xOy$ 中，椭圆的方程为 $\dfrac{x^2}{a^2}+\dfrac{y^2}{b^2}=1\ (a>b>0)$。$A_1,A_2$ 分别是椭圆的左、右顶点，$F_1,F_2$ 分别是椭圆的左、右焦点，$P$ 为椭圆上不同于 $A_1$ 和 $A_2$ 的任意一点。若平面中有两个点 $Q,R$ 满足 $QA_1\perp PA_1,\ QA_2\perp PA_2,\ RF_1\perp PF_1,\ RF_2\perp PF_2$，试确定线段 $QR$ 的长度与 $b$ 的大小关系，并给出证明。

\clearpage
\textbf{11.（本题满分20分）}设函数 $f(x)=ax^2+b$，求所有的正实数对 $(a,b)$，使得对任意实数 $x,y$，均有
$$
f(xy)+f(x+y)\ge f(x)f(y).
$$

\clearpage
\section*{2013年全国高中数学联合竞赛一试（B卷）}

一、填空题：本大题共8小题，每小题8分，满分64分。

1. 已知锐角三角形三条边的长度都是整数，且其中两条边的长度分别为 $3$ 和 $4$，则第三条边的长度为\blank{}。

2. 设 $i$ 是虚数单位，则
$$
i+2i^2+3i^3+\cdots+2013i^{2013}=\underline{\hspace{3em}}.
$$

3. 设集合 $A=\{2,0,1,3\}$，集合 $B=\{x\mid -x\in A,\ 2-x^2\notin A\}$，则集合 $B$ 中所有元素的和为\blank{}。

4. 已知正三棱锥 $P-ABC$ 的底面边长为 $1$，高为 $\sqrt2$，则其内切球的半径为\blank{}。

5. 在区间 $[0,\pi)$ 中，方程 $\sin 12x=x$ 的解的个数为\blank{}。

6. 定义在实数集 $\mathbb R$ 上的函数
$$
f(x)=\frac{\sin\pi x}{\sqrt{1+x+x^2}}
$$
的最小值为\blank{}。

7. 设 $a,b$ 为实数，函数 $f(x)=ax+b$ 满足：对任意 $x\in[0,1]$，有 $|f(x)|\le1$，则 $ab$ 的最大值为\blank{}。

8. 将正九边形的每个顶点等概率地涂上红、蓝两种颜色之一，则存在三个同色的顶点构成锐角三角形的概率为\blank{}。

二、解答题：本大题共3小题，满分56分。


\clearpage
\textbf{9.（本题满分16分）}已知数列 $\{a_n\}$ 满足 $a_1=2,\ a_n=2(n+a_{n-1}),\ n\ge2$。求数列 $\{a_n\}$ 的通项公式。

\clearpage
\textbf{10.（本题满分20分）}设 $a,b,c$ 是正实数，满足 $abc=1$。求证：
$$
a+b+c\le a^2+b^2+c^2.
$$

\clearpage
\textbf{11.（本题满分20分）}在平面直角坐标系 $xOy$ 中，点 $F$ 的坐标为 $(1,0)$，点 $A,B$ 在抛物线 $y^2=4x$ 上，满足
$$
\overrightarrow{OA}\cdot\overrightarrow{OB}=-4,\quad |\overrightarrow{FA}|-|\overrightarrow{FB}|=4\sqrt3.
$$
求 $\overrightarrow{FA}\cdot\overrightarrow{FB}$ 的值。

\clearpage
\section*{2011年全国高中数学联合竞赛一试（A卷）}

一、填空题：本大题共8小题，每小题8分，满分64分。

1. 设集合 $A=\{a_1,a_2,a_3,a_4\}$。若 $A$ 中所有三元子集的三个元素之和组成的集合为 $B=\{-1,3,5,8\}$，则集合 $A=$\blank{}。

2. 函数
$$
f(x)=\frac{\sqrt{x^2+1}}{x-1}
$$
的值域为\blank{}。

3. 设 $a,b$ 为正实数，且
$$
\frac1a+\frac1b\le2\sqrt2,\quad (a-b)^2=4(ab)^3,
$$
则 $\log_a b=$\blank{}。

4. 若
$$
\cos^5\theta-\sin^5\theta<7(\sin^3\theta-\cos^3\theta)\quad(\theta\in[0,2\pi)),
$$
则 $\theta$ 的取值范围为\blank{}。

5. 现安排 $7$ 名同学去参加 $5$ 个运动项目，要求甲、乙两同学不能参加同一个项目，每个项目都有人参加，每人只参加一个项目，则满足上述要求的不同安排方案数为\blank{}（用数字作答）。

6. 在四面体 $ABCD$ 中，已知
$$
\angle ADB=\angle BDC=\angle CDA=60^\circ,\quad AD=BD=3,\ CD=2.
$$
则四面体 $ABCD$ 的外接球的半径为\blank{}。

7. 直线 $x-2y-1=0$ 与抛物线 $y^2=4x$ 交于 $A,B$ 两点，$C$ 为抛物线上的一点，$\angle ACB=90^\circ$，则点 $C$ 的坐标为\blank{}。

8. 已知
$$
a_n=C_{200}^n(\sqrt[3]{6})^{200-n}\left(\frac1{\sqrt2}\right)^n\quad(n=1,2,\dots,95).
$$
则数列 $\{a_n\}$ 中整数项的个数为\blank{}。

二、解答题：本大题共3小题，满分56分。


\clearpage
\textbf{9.（本题满分16分）}设函数 $f(x)=|\lg(x+1)|$，实数 $a,b\ (a<b)$ 满足
$$
f(a)=f\left(-\frac{b+1}{b+2}\right),\quad f(10a+6b+21)=4\lg2.
$$
求 $a,b$ 的值。

\clearpage
\textbf{10.（本题满分20分）}已知数列 $\{a_n\}$ 满足：$a_1=2t-3\ (t\in\mathbb R,\ t\ne\pm1)$，
$$
a_{n+1}=\frac{(2t^{n+1}-3)a_n+2(t-1)t^n-1}{a_n+2t^n-1},\quad n\ge1.
$$
（1）求数列 $\{a_n\}$ 的通项公式；  
（2）若 $t>0$，试比较 $a_{n+1}$ 与 $a_n$ 的大小。

\clearpage
\textbf{11.（本题满分20分）}作斜率为 $\dfrac13$ 的直线 $l$ 与椭圆
$$
C:\frac{x^2}{36}+\frac{y^2}{4}=1
$$
交于 $A,B$ 两点，且 $P(3\sqrt2,\sqrt2)$ 在直线 $l$ 的左上方。  
（1）求证：$\triangle PAB$ 的内切圆的圆心在一条定直线上；  
（2）若 $\angle APB=60^\circ$，求 $\triangle PAB$ 的面积。

\clearpage
\section*{2011年全国高中数学联合竞赛一试（B卷）}

一、填空题：本大题共8小题，每小题8分，满分64分。

1. 设等差数列 $\{a_n\}$ 的前 $n$ 项和为 $S_n$，若 $S_{2010}-S_1=1$，则 $S_{2011}=$\blank{}。

2. 已知复数 $z$ 的模为 $1$，若 $z=z_1$ 和 $z=z_2$ 时 $|z+1+i|$ 分别取得最大值和最小值，则 $z_1-z_2=$\blank{}。

3. 设 $a,b$ 为正实数，且
$$
\frac1a+\frac1b\le2\sqrt2,\quad (a-b)^2=4(ab)^3,
$$
则 $\log_a b=$\blank{}。

4. 把扑克牌中的 $A,2,\dots,J,Q,K$ 分别看作数字 $1,2,\dots,11,12,13$。现将一副扑克牌中的黑桃、红桃各 $13$ 张放在一起，从中随机取出 $2$ 张牌，其花色相同且两个数的积是完全平方数的概率为\blank{}。

5. 若 $\triangle ABC$ 的角 $A,C$ 满足
$$
5(\cos A+\cos C)+4(\cos A\cos C+1)=0,
$$
则 $\tan\dfrac A2\tan\dfrac C2=$\blank{}。

6. 已知正方体 $ABCD-A_1B_1C_1D_1$ 的棱长为 $6$，$M,N$ 分别是 $BB_1,B_1C_1$ 上的点，满足 $B_1M=B_1N=2$。$S,P$ 分别是线段 $AD,MN$ 的中点，则异面直线 $SP$ 与 $AC_1$ 的距离为\blank{}。

7. 在 $\triangle ABC$ 中，$E,F$ 分别是 $AC,AB$ 的中点，$AB=\dfrac23 AC$。若 $\dfrac{BE}{CF}<t$ 恒成立，则 $t$ 的最小值为\blank{}。

8. 抛物线
$$
y^2=2p\left(x-\frac p2\right)\quad(p>0)
$$
上动点 $A$ 到点 $B(3,0)$ 的距离的最小值记为 $d(p)$。则满足 $d(p)=2$ 的所有实数 $p$ 的和为\blank{}。

二、解答题：本大题共3小题，满分56分。


\clearpage
\textbf{9.（本题满分16分）}已知实数 $x,y,z$ 满足
$$
x\ge y\ge z,\quad x+y+z=1,\quad x^2+y^2+z^2=3.
$$
求 $x$ 的取值范围。

\clearpage
\textbf{10.（本题满分20分）}已知数列 $\{a_n\}$ 满足：$a_1=2t-2\ (t\in\mathbb R,\ t\ne\pm1)$，
$$
a_{n+1}=\frac{2(t^{n+1}-1)a_n}{a_n+2t^n-2},\quad n\ge1.
$$
（1）求数列 $\{a_n\}$ 的通项公式；  
（2）若 $t>0$，试比较 $a_{n+1}$ 与 $a_n$ 的大小。

\clearpage
\textbf{11.（本题满分20分）}已知 $A_1(x_1,y_1),A_2(x_2,y_2),A_3(x_3,y_3)$ 是抛物线 $y^2=2px\ (p>0)$ 上不同的三点，$\triangle A_1A_2A_3$ 有两边所在的直线与抛物线 $x^2=2qy\ (q>0)$ 相切。求证：对不同的 $i,j\in\{1,2,3\}$，$y_iy_j(y_i+y_j)$ 为定值。

\clearpage
\section*{2012年全国高中数学联合竞赛一试（A卷）}

一、填空题：本大题共8小题，每小题8分，满分64分。

1. 设 $P$ 是函数 $y=x+\dfrac2x\ (x>0)$ 的图象上的任意一点，过 $P$ 分别向直线 $y=x$ 和 $y$ 轴作垂线，垂足分别为 $A,B$。则 $\overrightarrow{PA}\cdot\overrightarrow{PB}$ 的值为\blank{}。

2. 设 $\triangle ABC$ 的内角 $\angle A,\angle B,\angle C$ 的对边分别为 $a,b,c$，且满足
$$
a\cos B-b\cos A=\frac35 c.
$$
则 $\dfrac{\tan A}{\tan B}$ 的值为\blank{}。

3. 设 $x,y,z\in[0,1]$，则
$$
M=\sqrt{|x-y|}+\sqrt{|y-z|}+\sqrt{|z-x|}
$$
的最大值为\blank{}。

4. 设抛物线 $y^2=2px\ (p>0)$ 的焦点为 $F$，准线为 $l$，$A,B$ 是抛物线上的两个动点，且满足 $\angle AFB=\dfrac\pi3$。设线段 $AB$ 的中点 $M$ 在 $l$ 上的投影为 $N$，则 $\dfrac{|MN|}{|AB|}$ 的最大值为\blank{}。

5. 设同底的两个正三棱锥 $P-ABC$ 和 $Q-ABC$ 内接于同一个球。若正三棱锥 $P-ABC$ 的侧面与底面所成的角为 $45^\circ$，则正三棱锥 $Q-ABC$ 的侧面与底面所成角的正切值为\blank{}。

6. 设 $f(x)$ 是定义在 $\mathbb R$ 上的奇函数，且当 $x\ge0$ 时，$f(x)=x^2$。若对任意的 $x\in[a,a+2]$，不等式 $f(x+a)\ge2f(x)$ 恒成立，则实数 $a$ 的取值范围为\blank{}。

7. 满足
$$
\frac14<\sin\frac\pi n<\frac13
$$
的所有正整数 $n$ 的和为\blank{}。

8. 某情报站有 $A,B,C,D$ 四种互不相同的密码，每周使用其中的一种密码，且每周都是从上周末使用的三种密码中等可能地随机选用一种。设第一周使用 $A$ 种密码，那么第七周也使用 $A$ 种密码的概率为\blank{}。

二、解答题：本大题共3小题，满分56分。


\clearpage
\textbf{9.（本题满分16分）}已知函数
$$
f(x)=a\sin x-\frac12\cos2x+a-\frac3a+\frac12,
$$
其中 $a\in\mathbb R$，且 $a\ne0$。  
（1）若对任意 $x\in\mathbb R$，都有 $f(x)\le0$，求 $a$ 的取值范围；  
（2）若 $a\ge2$，且存在 $x\in\mathbb R$，使 $f(x)\le0$，求 $a$ 的取值范围。

\clearpage
\textbf{10.（本题满分20分）}已知数列 $\{a_n\}$ 的各项均为非零实数，且对任意的正整数 $n$，都有
$$
(a_1+a_2+\cdots+a_n)^2=a_1^3+a_2^3+\cdots+a_n^3.
$$
（1）当 $n=3$ 时，求所有满足条件的三项组成的数列 $a_1,a_2,a_3$；  
（2）是否存在满足条件的无穷数列 $\{a_n\}$，使得 $a_{2013}=-2012$？若存在，求出这样的无穷数列的一个通项公式；若不存在，说明理由。

\clearpage
\textbf{11.（本题满分20分）}在平面直角坐标系 $xOy$ 中，菱形 $ABCD$ 的边长为 $4$，且 $|OB|=|OD|=6$。  
（1）求证：$|OA|\cdot|OC|$ 为定值；  
（2）当点 $A$ 在半圆
$$
M:(x-2)^2+y^2=4\quad(2\le x\le4)
$$
上运动时，求点 $C$ 的轨迹。

\clearpage
\section*{2012年全国高中数学联合竞赛一试（B卷）}

一、填空题：本大题共8小题，每小题8分，满分64分。

1. 对于集合 $\{x\mid a\le x\le b\}$，我们把 $b-a$ 称为它的长度。设集合
$$
A=\{x\mid a\le x\le a+1981\},\quad B=\{x\mid b-1014\le x\le b\},
$$
且 $A,B$ 都是集合 $U=\{x\mid0\le x\le2012\}$ 的子集。则集合 $A\cap B$ 的长度的最小值为\blank{}。

2. 已知 $x>0,\ y>0$，且满足
$$
\begin{cases}
\cos^2\pi x+2\sin\pi y=1,\\
\sin\pi x+\sin\pi y=0,\\
x^2-y^2=12,
\end{cases}
$$
则有序数对 $(x,y)=$\blank{}。

3. 设椭圆 $\dfrac{x^2}{a^2}+\dfrac{y^2}{b^2}=1\ (a>b>0)$ 的左、右焦点分别为 $F_1,F_2$，过点 $F_2$ 的直线交椭圆于 $A(x_1,y_1),B(x_2,y_2)$ 两点。若 $\triangle AF_1B$ 内切圆的面积为 $\pi$，且 $|y_1-y_2|=4$，则椭圆的离心率为\blank{}。

4. 若关于 $x$ 的不等式组
$$
\begin{cases}
x^3+3x^2-x-3>0,\\
x^2-2ax-1\le0
\end{cases}\quad(a>0)
$$
的整数解有且只有一个，则 $a$ 的取值范围为\blank{}。

5. 在 $\triangle ABC$ 中，若
$$
\overrightarrow{AB}\cdot\overrightarrow{AC}=7,\quad |\overrightarrow{AB}-\overrightarrow{AC}|=6,
$$
则 $\triangle ABC$ 面积的最大值为\blank{}。

6. 在长方体 $ABCD-A_1B_1C_1D_1$ 中，$AB=4,\ BC=CC_1=2\sqrt2$，$M$ 是 $BC_1$ 的中点，$N$ 是 $MC_1$ 的中点。若异面直线 $AN$ 与 $CM$ 所成的角为 $\theta$、距离为 $d$，则 $d\sin\theta=$\blank{}。

7. 设 $f(x)$ 是定义在 $\mathbb R$ 上的奇函数，且当 $x\ge0$ 时，$f(x)=x^2$。若对任意的 $x\in[a,a+2]$，不等式 $f(x+a)\ge2f(x)$ 恒成立，则实数 $a$ 的取值范围为\blank{}。

8. 一个均匀的正方体骰子的各面上分别标有数字 $1,2,\dots,6$。每次投掷两个相同的骰子，规定向上的两个面上的数字之和为这次投掷的点数。那么，投掷 $3$ 次所得 $3$ 个点数之积能被 $14$ 整除的概率为\blank{}。

二、解答题：本大题共3小题，满分56分。


\clearpage
\textbf{9.（本题满分16分）}已知函数
$$
f(x)=a\sin x-\frac12\cos2x+a-\frac3a+\frac12,
$$
其中 $a\in\mathbb R$，且 $a\ne0$。  
（1）若对任意 $x\in\mathbb R$，都有 $f(x)\le0$，求 $a$ 的取值范围；  
（2）若 $a\ge2$，且存在 $x\in\mathbb R$，使 $f(x)\le0$，求 $a$ 的取值范围。

\clearpage
\textbf{10.（本题满分20分）}已知数列 $\{a_n\}$ 满足：当 $n=2k-1\ (k\in\mathbb N^\emph{)$ 时，$a_n=n$；当 $n=2k\ (k\in\mathbb N^})$ 时，$a_n=a_k$。记
$$
T_n=a_1+a_2+a_3+\cdots+a_{2^n}.
$$
求证：对任意正整数 $n$，有  
（1）$T_{n+1}=4^n+T_n$；  
（2）
$$
\frac1{T_1}+\frac1{T_2}+\cdots+\frac1{T_n}<1.
$$

\clearpage
\textbf{11.（本题满分20分）}已知抛物线 $y^2=2px\ (p>0)$，$A,B$ 是抛物线上不同于顶点 $O$ 的两个动点。记 $\angle AOB=\theta\ (\theta\ne90^\circ)$，若 $S_{\triangle AOB}=m\tan\theta$，试求当 $m$ 取到最小值时 $\tan\theta$ 的最大值。


以下是为您提取的四道大题的文字内容（已按照题号顺序排版，公式使用 LaTeX 格式）：

\clearpage
\textbf{第9题}
(2018)设 $n$ 为正整数，集合
$$
A = \{\alpha \mid \alpha = (t_1, t_2, \cdots, t_n), t_k \in \{0, 1\}, k = 1, 2, \cdots, n\}.
$$

对于集合 $A$ 中的任意元素 $\alpha = (x_1, x_2, \cdots, x_n)$ 和 $\beta = (y_1, y_2, \cdots, y_n)$，记
$$
M(\alpha, \beta) = \frac{1}{2}[(x_1 + y_1 - |x_1 - y_1|) + \cdots + (x_n + y_n - |x_n - y_n|)].
$$

(1) 当 $n = 3$ 时，若 $\alpha = (1, 1, 0), \beta = (0, 1, 1)$，求 $M(\alpha, \alpha)$ 和 $M(\alpha, \beta)$ 的值；

(2) 当 $n = 4$ 时，设 $B$ 是 $A$ 的子集，且满足：对于 $B$ 中的任意元素 $\alpha, \beta$，当 $\alpha, \beta$ 相同时，$M(\alpha, \beta)$ 是奇数；当 $\alpha, \beta$ 不同时，$M(\alpha, \beta)$ 是偶数. 求集合 $B$ 中元素个数的最大值；

(3) 给定不小于 $2$ 的 $n$，设 $B$ 是 $A$ 的子集，且满足：对于 $B$ 中任意两个不同的元素 $\alpha, \beta$，$M(\alpha, \beta) = 0$. 写出一个集合 $B$，使其元素个数最多，并说明理由.

\clearpage
\textbf{第10题}
(2017)设 $\{a_n\}$ 和 $\{b_n\}$ 是两个等差数列，记
$$
c_n = \max\{b_1 - a_1n, b_2 - a_2n, \cdots, b_n - a_nn\} (n = 1, 2, 3, \cdots),
$$
其中 $\max\{x_1, x_2, \cdots, x_s\}$ 表示 $x_1, x_2, \cdots, x_s$ 这 $s$ 个数中最大的数.

(1) 若 $a_n = n, b_n = 2n - 1$，求 $c_1, c_2, c_3$ 的值，并证明 $\{c_n\}$ 是等差数列；

(2) 证明：或者对于任意正数 $M$，存在正整数 $m$，当 $n \ge m$ 时，$\frac{c_n}{n} > M$，或者存在正整数 $m$，使得 $c_m, c_{m+1}, c_{m+2}, \cdots$ 是等差数列.

\clearpage
\textbf{第11题}
(2016)设数列 $A : a_1, a_2, \cdots, a_N (N \ge 2)$，如果对于小于 $n (2 \le n \le N)$ 的每个正整数 $k$ 都有 $a_k < a_n$，则称 $n$ 是数列 $A$ 的一个“$G$ 时刻”. 记 $G(A)$ 是数列 $A$ 的所有“$G$ 时刻”组成的集合.

(1) 对数列 $A : -2, 2, -1, 1, 3$，写出 $G(A)$ 的所有元素；

(2) 证明：若数列 $A$ 中存在 $a_n > a_1$，则 $G(A) \neq \varnothing$；

(3) 证明：若数列 $A$ 满足 $a_n - a_{n-1} \le 1 (n = 2, 3, \cdots, N)$，则 $G(A)$ 的元素个数不小于 $a_N - a_1$.

\clearpage
\textbf{第12题}
(2015)已知数列 $\{a_n\}$ 满足：$a_1 \in \mathbb{N}^*, a_1 \le 36$，且
$$
a_{n+1} = \begin{cases} 2a_n, & a_n \le 18 \\ 2a_n - 36, & a_n > 18 \end{cases} (n = 1, 2, \cdots).
$$
记集合 $M = \{a_n \mid n \in \mathbb{N}^*\}$.

(1) 若 $a_1 = 6$，写出集合 $M$ 的所有元素；

(2) 若集合 $M$ 中存在一个元素是 3 的倍数，证明：$M$ 的所有元素都是 3 的倍数；

(3) 求集合 $M$ 的元素个数的最大值.

\end{document}